
\paragraph{\SGARefLink{sga3:VIIB:5.2.1.B}{5.2.1.B}.}
\phantomsection
\label{sga3:VIIB:5.2.1.B}

As above, put
\[
  \omega_\infty
  =
  \overline{\bigcup_{n\in\mathbb N}\omega_n}.
\]
A second case in which a section
\[
  \sigma:\omega\longrightarrow I
\]
of the desired kind exists is the case in which
\(\omega_\infty\) has a pseudobasis \(\mathscr B_\infty\) that
is the union, for \(n\in\mathbb N^*\), of pseudobases of the
spaces
\[
  \omega_n/\omega_{n-1}.
\]
One can then complete it by a pseudobasis
\(\mathscr B'_\infty\) of \(\omega/\omega_\infty\), obtaining a
pseudobasis of \(\omega\) compatible with the filtration.

Put
\[
  V=\omega_\infty^\dagger,
\]
and let \(V_n\) be the orthogonal complement of \(\omega_n\) in
\(V\). The preceding condition is equivalent to saying that, in
the category of ordinary \(k\)-vector spaces, the separated
decreasing filtration
\[
  V=V_0\supset V_1\supset V_2\supset\cdots
\]
splits: namely, \(V\) is the direct sum, over
\(n\in\mathbb N\), of subspaces \(F_n\) such that
\[
  F_n\simeq V_n/V_{n+1}.
\]

This need not hold. For example, let
\[
  V=\mathscr S=k^{\mathbb N}
\]
be the space of sequences of elements of \(k\), and let
\(\mathscr S_n\) be the subspace of sequences
\((u_i)\) such that \(u_i=0\) for \(i<n\). Then
\[
  \dim\mathscr S_n/\mathscr S_{n+1}=1,
\]
but \(\mathscr S\) is not isomorphic to the direct sum of the
\(\mathscr S_n/\mathscr S_{n+1}\), because \(\mathscr S\) is not
countable-dimensional. In this example, however, \(\mathscr S\)
is the product of the \(\mathscr S_n/\mathscr S_{n+1}\)
\textup{(cf.\ \SGARefLink{sga3:VIIB:5.2.1.A}{5.2.1.A})}.

The splitting does hold when \(V\) is
countable-dimensional.
Indeed, let
\[
  (e_n)_{n\in\mathbb N}
\]
be a basis of \(V\). We shall construct inductively an increasing
function
\[
  g:\mathbb N\longrightarrow\mathbb N
\]
and subspaces \(F_i\), for \(i=0,\ldots,g(n)\), such that
\[
  F_i\simeq V_i/V_{i+1}
\]
and
\[
  F_{\leq g(n)}
  =
  \bigoplus_{i=0}^{g(n)}F_i
\]
is a complement of \(V_{g(n)+1}\) containing
\[
  e_0,\ldots,e_n.
\]
It will then follow that
\[
  V=\bigoplus_{i\geq0}F_i.
\]

Let \(n+1\in\mathbb N\), and assume the assertion already proved
for \(n\) \textup{(it is vacuous for \(n=-1\))}. If
\[
  e_{n+1}\in F_{\leq g(n)},
\]
put \(g(n+1)=g(n)\). Otherwise write
\[
  e_{n+1}=f+x,
  \qquad
  f\in F_{\leq g(n)},
  \qquad
  0\ne x\in V_{g(n)+1}.
\]
Let \(j\) be the least integer such that
\[
  x\in V_j\setminus V_{j+1}.
\]
For
\[
  i=g(n)+1,\ldots,j,
\]
choose a complement \(F_i\) of \(V_{i+1}\) in \(V_i\), choosing
\(F_j\) so that \(x\in F_j\). Then
\[
  e_{n+1}\in F_{\leq j};
\]
put \(g(n+1)=j\).

\phantomsection
\label{sga3:VIIB:5.2.1.B:induction-source-note}
\noindent\textit{Translator's note.}\enspace At the last induction step, the French text prints
\(e_{n+1}\in F_j\). From
\(e_{n+1}=f+x\), with \(f\in F_{\leq g(n)}\), the argument
instead gives \(x\in F_j\) and
\(e_{n+1}\in F_{\leq j}\); these are the relations used here.
